\subsection{Phương Pháp Ordinary Least Squares (OLS)}


\subsubsection{Hàm Mất Mát và Nghiệm OLS}

OLS tìm $\hat{\boldsymbol{\beta}}$ tối thiểu hóa \emph{Tổng Bình Phương Phần
Dư} (Residual Sum of Squares - RSS):
\begin{equation}
    \mathrm{RSS}(\boldsymbol{\beta})
    = \|\mathbf{y} - \mathbf{X}\boldsymbol{\beta}\|_2^2
    = \sum_{i=1}^n (y_i - \mathbf{x}_i^T\boldsymbol{\beta})^2.
    \label{eq:rss}
\end{equation}

\begin{tcolorbox}[mycodebox, colback=blue!3!white, colframe=blue!40!black,
    title={\textbf{Định lý 1.1} - Nghiệm OLS (Normal Equations)}]
Nếu $\mathbf{X}^T\mathbf{X}$ khả nghịch, nghiệm OLS duy nhất là:
\begin{equation}
    \hat{\boldsymbol{\beta}}_{\mathrm{OLS}}
    = (\mathbf{X}^T\mathbf{X})^{-1}\mathbf{X}^T\mathbf{y}.
    \label{eq:ols_solution}
\end{equation}
\end{tcolorbox}

\textbf{Chứng minh.} Lấy đạo hàm của RSS theo $\boldsymbol{\beta}$:
\begin{align}
    \nabla_{\boldsymbol{\beta}}\,\mathrm{RSS}
    &= \nabla_{\boldsymbol{\beta}}
       \bigl[(\mathbf{y} - \mathbf{X}\boldsymbol{\beta})^T
             (\mathbf{y} - \mathbf{X}\boldsymbol{\beta})\bigr] \notag \\
    &= -2\,\mathbf{X}^T(\mathbf{y} - \mathbf{X}\boldsymbol{\beta}).
\end{align}
Cho đạo hàm bằng $\mathbf{0}$:
$\mathbf{X}^T\mathbf{X}\boldsymbol{\beta} = \mathbf{X}^T\mathbf{y}$
(phương trình pháp tuyến - \emph{Normal Equations}).
Vì $\mathbf{X}^T\mathbf{X}$ khả nghịch (theo GM2), suy ra
$\hat{\boldsymbol{\beta}} = (\mathbf{X}^T\mathbf{X})^{-1}\mathbf{X}^T\mathbf{y}$.
Để kiểm tra đây là cực tiểu (không phải cực đại), ta kiểm tra đạo hàm bậc 2:
$\nabla^2_{\boldsymbol{\beta}}\,\mathrm{RSS} = 2\mathbf{X}^T\mathbf{X}$
là ma trận xác định dương (positive definite) vì $\rank(\mathbf{X}) = p+1$.
$\hfill\square$

\textbf{Cài đặt Python - F1: \texttt{ols\_fit}:}

\begin{lstlisting}[language=Python, caption={Hàm \texttt{ols\_fit} - tính nghiệm OLS từ công thức Normal Equations}]
def ols_fit(X: list[list[float]], y: list[float]) -> dict:
    """F1: Tinh beta_hat = (X^T X)^{-1} X^T y va sigma2_hat."""
    n, p = len(X), len(X[0])
    X_b  = add_bias(X)                        # them cot 1 (intercept)
    Xt   = transpose(X_b)                     # utils.py
    XtX  = matmul(Xt, X_b)                    # utils.py
    Xty  = matvec(Xt, y)                      # utils.py
    beta_hat  = solve_system(XtX, Xty)        # utils.py (Gauss elim.)
    y_hat     = matvec(X_b, beta_hat)
    residuals = [y[i] - y_hat[i] for i in range(n)]
    rss       = sum(r * r for r in residuals)
    sigma2    = rss / max(n - p - 1, 1)       # ung luong khong chech
    return {"beta_hat": beta_hat, "sigma2_hat": sigma2,
            "y_hat": y_hat, "residuals": residuals}
\end{lstlisting}

Hàm \texttt{solve\_system} trong \texttt{utils.py} giải hệ $\mathbf{A}\mathbf{x}
= \mathbf{b}$ bằng phép khử Gauss với partial pivoting, tránh sử dụng
\texttt{numpy.linalg.solve} đúng theo yêu cầu đề bài.

\subsubsection{Ma Trận Chiếu và Hat Matrix}

\begin{tcolorbox}[mycodebox, colback=blue!3!white, colframe=blue!40!black,
    title={\textbf{Định nghĩa 1.2} - Hat Matrix}]
Ma trận chiếu (projection matrix) hay \emph{Hat Matrix} được định nghĩa:
\begin{equation}
    \mathbf{H} = \mathbf{X}(\mathbf{X}^T\mathbf{X})^{-1}\mathbf{X}^T
    \in \mathbb{R}^{n \times n}.
    \label{eq:hat_matrix}
\end{equation}
\end{tcolorbox}

Tên ``Hat Matrix'' xuất phát từ tính chất $\hat{\mathbf{y}} = \mathbf{H}\mathbf{y}$
- ma trận ``đội mũ'' lên $\mathbf{y}$ để tạo ra giá trị dự đoán $\hat{\mathbf{y}}$.

\begin{tcolorbox}[mycodebox, colback=green!3!white, colframe=green!40!black,
    title={\textbf{Mệnh đề 1.1} - Tính Chất của $\mathbf{H}$}]
\begin{enumerate}[label=(\roman*)]
    \item $\mathbf{H}^2 = \mathbf{H}$ \hfill (idempotent - lũy đẳng)
    \item $\mathbf{H}^T = \mathbf{H}$ \hfill (đối xứng)
    \item Giá trị riêng của $\mathbf{H}$: chỉ là $0$ hoặc $1$
    \item $\rank(\mathbf{H}) = \mathrm{tr}(\mathbf{H}) = p+1$
    \item $\hat{\mathbf{y}} = \mathbf{H}\mathbf{y}$;\quad
          $\hat{\boldsymbol{\varepsilon}} = (\mathbf{I} - \mathbf{H})\mathbf{y}$
    \item $(\mathbf{I}-\mathbf{H})$ cũng idempotent và đối xứng
\end{enumerate}
\end{tcolorbox}

\textbf{Chứng minh các tính chất.}
\begin{itemize}
    \item \textit{(i) Idempotent:}
    $\mathbf{H}^2 = \mathbf{X}(\mathbf{X}^T\mathbf{X})^{-1}\mathbf{X}^T
                   \mathbf{X}(\mathbf{X}^T\mathbf{X})^{-1}\mathbf{X}^T
               = \mathbf{X}(\mathbf{X}^T\mathbf{X})^{-1}\mathbf{X}^T = \mathbf{H}$.
    \item \textit{(ii) Đối xứng:}
    $\mathbf{H}^T = \bigl[\mathbf{X}(\mathbf{X}^T\mathbf{X})^{-1}\mathbf{X}^T\bigr]^T
               = \mathbf{X}\bigl[(\mathbf{X}^T\mathbf{X})^{-1}\bigr]^T\mathbf{X}^T
               = \mathbf{X}(\mathbf{X}^T\mathbf{X})^{-1}\mathbf{X}^T = \mathbf{H}$
    (vì $(\mathbf{X}^T\mathbf{X})$ đối xứng nên nghịch đảo cũng đối xứng).
    \item \textit{(iii) Giá trị riêng:} Nếu $\lambda$ là giá trị riêng của $\mathbf{H}$,
    từ $\mathbf{H}^2 = \mathbf{H}$ suy ra $\lambda^2 = \lambda$, tức $\lambda \in \{0, 1\}$.
    \item \textit{(iv) Rank:} Vì $\mathbf{H}$ idempotent,
    $\rank(\mathbf{H}) = \mathrm{tr}(\mathbf{H})
    = \mathrm{tr}\bigl[\mathbf{X}(\mathbf{X}^T\mathbf{X})^{-1}\mathbf{X}^T\bigr]
    = \mathrm{tr}\bigl[(\mathbf{X}^T\mathbf{X})^{-1}\mathbf{X}^T\mathbf{X}\bigr]
    = \mathrm{tr}(\mathbf{I}_{p+1}) = p+1$.
\end{itemize}

\subsubsection{Định Lý Gauss-Markov}

\begin{tcolorbox}[mycodebox, colback=blue!3!white, colframe=blue!40!black,
    title={\textbf{Định lý 1.2} - Gauss-Markov (BLUE)}]
Dưới các giả thiết GM1-GM4, ước lượng OLS
$\hat{\boldsymbol{\beta}}_{\mathrm{OLS}}$ là ước lượng tuyến tính không
chệch tốt nhất (\emph{Best Linear Unbiased Estimator} - \textbf{BLUE}):
\begin{enumerate}[label=(\roman*)]
    \item \textbf{Không chệch (Unbiased):}
          $E[\hat{\boldsymbol{\beta}}_{\mathrm{OLS}}] = \boldsymbol{\beta}$.
    \item \textbf{Tốt nhất (Best):} Với mọi ước lượng tuyến tính không chệch
          $\tilde{\boldsymbol{\beta}}$ khác,
          $\mathrm{Var}(\tilde{\beta}_j) \geq \mathrm{Var}(\hat{\beta}_j^{\mathrm{OLS}})$
          với mọi $j$.
\end{enumerate}
Ma trận hiệp phương sai của $\hat{\boldsymbol{\beta}}_{\mathrm{OLS}}$:
\begin{equation}
    \mathrm{Var}(\hat{\boldsymbol{\beta}}_{\mathrm{OLS}} \mid \mathbf{X})
    = \sigma^2(\mathbf{X}^T\mathbf{X})^{-1}.
    \label{eq:var_beta}
\end{equation}
\end{tcolorbox}

\textbf{Chứng minh tính không chệch.}
\begin{align}
    E[\hat{\boldsymbol{\beta}}_{\mathrm{OLS}}]
    &= E\bigl[(\mathbf{X}^T\mathbf{X})^{-1}\mathbf{X}^T\mathbf{y}\bigr] \notag \\
    &= (\mathbf{X}^T\mathbf{X})^{-1}\mathbf{X}^T
       E[\mathbf{X}\boldsymbol{\beta} + \boldsymbol{\varepsilon}] \notag \\
    &= (\mathbf{X}^T\mathbf{X})^{-1}\mathbf{X}^T\mathbf{X}\boldsymbol{\beta}
       + (\mathbf{X}^T\mathbf{X})^{-1}\mathbf{X}^T
         \underbrace{E[\boldsymbol{\varepsilon} \mid \mathbf{X}]}_{=\mathbf{0}\text{ (GM3)}}
    = \boldsymbol{\beta}. \quad \square
\end{align}

\textbf{Chứng minh ma trận hiệp phương sai.}
\begin{align}
    \mathrm{Var}(\hat{\boldsymbol{\beta}})
    &= \mathrm{Var}\bigl[(\mathbf{X}^T\mathbf{X})^{-1}\mathbf{X}^T\mathbf{y}\bigr]
     = (\mathbf{X}^T\mathbf{X})^{-1}\mathbf{X}^T
       \underbrace{\mathrm{Var}(\mathbf{y})}_{=\sigma^2\mathbf{I}_n}
       \mathbf{X}(\mathbf{X}^T\mathbf{X})^{-1}
     = \sigma^2(\mathbf{X}^T\mathbf{X})^{-1}. \quad \square
\end{align}

\textbf{Minh họa Monte Carlo - Định lý Gauss-Markov (F10):}

Để kiểm chứng định lý bằng mô phỏng, ta thực hiện $N_{\mathrm{sim}} = 1000$
lần với cùng ma trận $\mathbf{X}$ cố định, mỗi lần sinh ngẫu nhiên
$\boldsymbol{\varepsilon} \sim \mathcal{N}(\mathbf{0}, \sigma^2\mathbf{I})$
và tính $\hat{\boldsymbol{\beta}}_{\mathrm{OLS}}$. Bên cạnh đó, xây dựng
estimator thay thế $\tilde{\boldsymbol{\beta}} = \hat{\boldsymbol{\beta}} +
c\,\mathbf{e}_0 (\mathbf{v}^T\mathbf{y})$ với
$\mathbf{v} \in \mathrm{Null}(\mathbf{X}^T)$
(vẫn là ước lượng không chệch nhưng có phương sai lớn hơn).

Trước khi quay lại phần cài đặt, Hình~\ref{fig:gauss_markov} đóng vai trò như một kiểm chứng
thực nghiệm cho phần chứng minh ở trên: nếu các giả định Gauss--Markov được giữ, OLS không chỉ
không chệch mà còn có phương sai nhỏ nhất trong lớp ước lượng tuyến tính không chệch được thử.
Kết quả mô phỏng xác nhận:
\begin{itemize}
    \item $E[\hat{\boldsymbol{\beta}}_{\mathrm{OLS}}] \approx \boldsymbol{\beta}$
          (sai số trung bình $< 0.01$).
    \item $\mathrm{Var}(\hat{\beta}_j^{\mathrm{OLS}}) \leq
          \mathrm{Var}(\tilde{\beta}_j)$ với mọi $j$.
\end{itemize}

\begin{figure}[h]
    \centering
    \includegraphics[width=0.92\textwidth]{assets/part1/gauss_markov_histograms.png}
    \caption{Phân phối $\hat{\boldsymbol{\beta}}_{\mathrm{OLS}}$ (xanh) và
             $\tilde{\boldsymbol{\beta}}_{\mathrm{alt}}$ (cam) qua 1000 mô phỏng.
             Đường đỏ dọc là giá trị $\beta$ thực.
             OLS tập trung xung quanh $\beta$ thực (không chệch) với độ phân
             tán thấp hơn estimator thay thế (phương sai nhỏ hơn).}
    \label{fig:gauss_markov}
\end{figure}

\subsubsection{Ước Lượng Phương Sai Nhiễu}

\begin{tcolorbox}[mycodebox, colback=green!3!white, colframe=green!40!black,
    title={\textbf{Định lý 1.3} - Ước Lượng Không Chệch của $\sigma^2$}]
\begin{equation}
    \hat{\sigma}^2
    = \frac{\mathrm{RSS}}{n - p - 1}
    = \frac{\|\mathbf{y} - \mathbf{X}\hat{\boldsymbol{\beta}}\|^2}{n - p - 1},
    \label{eq:sigma2}
\end{equation}
với $n - p - 1$ là bậc tự do còn lại sau khi ước lượng $p+1$ tham số.
Ta có $E[\hat{\sigma}^2] = \sigma^2$.
\end{tcolorbox}

\textbf{Lưu ý triển khai:} Hàm \texttt{ols\_fit} tính $\hat{\sigma}^2$
sử dụng \texttt{n - p - 1} bậc tự do (không phải $n - 2$),
đảm bảo tính đúng khi số biến $p > 1$.


